\documentclass[../../../include/open-logic-section]{subfiles}

\begin{document}

\olfileid[pt]{sth}{z}{power}
\olsection{Conjuntos potência}

Prosseguiremos com outro axioma:

\begin{axiom}[Conjuntos potência]
Para qualquer conjunto $A$, existe o conjunto $\Pow{A} = \Setabs{x}{x \subseteq A}$.
\[
	\forall A \exists P \forall x(x \in P \liff (\forall z \in x)z \in A)
\]
\end{axiom}

A nossa justificação é bastante direta. Suponhamos que $A$ se forma no estágio $S$. Então todos os membros de $A$ estavam disponíveis antes de
$S$ (por \stagesacc). Logo, raciocinando como na justificação da
Separação, cada subconjunto de $A$ forma-se até ao estágio $S$. Estão
todos disponíveis, para serem reunidos num único conjunto, em qualquer estágio depois de $S$.
E sabemos que existe algum tal estágio, pois $S$ não é o último
estágio (por \stagessucc). Logo $\Pow{A}$ existe.

Eis uma boa consequência dos conjuntos potência:

\begin{prop}\label{thm:Products}
Dados quaisquer conjuntos $A, B$, o seu produto cartesiano $A \times B$ existe.
\end{prop}

\begin{proof}
O conjunto $\Pow{\Pow{A \cup B}}$ existe por Conjuntos potência e
\olref[sth][z][pairs]{prop:pairsconsequences}. Logo, por Separação, existe este
conjunto:
\[
	C = \Setabs{z \in \Pow{\Pow{A \cup B}}}{(\exists x \in A)(\exists y \in B) z = \tuple{x, y}}.
\]
Agora, para quaisquer $x \in A$ e $y \in B$, o conjunto $\tuple{x, y}$ existe por
\olref[sth][z][pairs]{prop:pairsconsequences}. Além disso, uma vez que $x, y
\in A \cup B$, temos $\{x\}, \{x, y\} \in \Pow{A \cup B}$, e
$\tuple{x,y} \in \Pow{\Pow{A \cup B}}$. Logo $A \times B = C$.
\end{proof}

Nesta prova, o axioma dos conjuntos potência interage com a Separação. E isso não
surpreende. Sem Separação, o princípio dos conjuntos potência não seria muito
\emph{poderoso}. Afinal, a Separação nos diz que subconjuntos de um conjunto existem e, por conseguinte, determina quão ``gordo'' é cada
conjunto potência.

\begin{prob}
Mostre que, quaisquer que sejam os conjuntos $A, B$: (i) existe o conjunto de todas as relações com
domínio $A$ e contradomínio $B$; e (ii) existe o conjunto de todas as funções
de $A$ para $B$.
\end{prob}

\begin{prob}
Seja $A$ um conjunto, e seja $\sim$ uma relação de equivalência em $A$.
Prove que existe o conjunto das classes de equivalência sob $\sim$ em $A$, isto é,
$\equivclass{A}{\sim}$.
\end{prob}

\end{document}
